\begin{solution}{10}
如图，在小球 $1$ 未离开竖直墙面之前，杆与竖直墙面之间夹角为 $\theta$ 时，小球 $1$ 的坐标为
\begin{equation}
x _ {1} = 0, y _ {1} = l \cos \theta ,
\label{eq:1.s10}
\end{equation}
小球 $2$ 的坐标为
\begin{equation}
x _ {2} = l \sin \theta , y _ {2} = 0
\label{eq:2.s10}
\end{equation}
小球 $1$ 的速度为
\begin{equation}
v _ {1 x} = \frac {\mathrm{d} x _ {1}}{\mathrm{d} t} = 0, v _ {1 y} = \frac {\mathrm{d} y _ {1}}{\mathrm{d} t} = - \omega l \sin \theta ,
\label{eq:3.s10}
\end{equation}
式中 $\omega=\frac{d\theta}{dt}$ 是杆的转动角速度。小球 $2$ 的速度为
\begin{equation}
v _ {2 x} = \frac {\mathrm{d} x _ {2}}{\mathrm{d} t} = \omega l \cos \theta , v _ {2 y} = \frac {\mathrm{d} y _ {2}}{\mathrm{d} t} = 0
\label{eq:4.s10}
\end{equation}
由机械能守恒，有
\begin{equation}
m g l = \frac {1}{2} m \left(v _ {1 x} ^ {2} + v _ {1 y} ^ {2}\right) + \frac {1}{2} m \left(v _ {2 x} ^ {2} + v _ {2 y} ^ {2}\right) + m g l \cos \theta = \frac {1}{2} m l ^ {2} \omega^ {2} + m g l \cos \theta 
\label{eq:5.s10}
\end{equation}
由上式得
\begin{equation}
\omega = \sqrt {\frac {2 g (1 - \cos \theta)}{l}}
\label{eq:6.s10}
\end{equation}
这里考虑到随着时间 $t$ 的增加， $\theta$ 变大，因此 $\omega > 0$ ，从而舍去负根。

系统质心 $c$ 的 $x$ 坐标为
\begin{equation}
x _ {c} = \frac {m x _ {1} + m x _ {2}}{2 m} = \frac {1}{2} l \sin \theta 
\label{eq:7.s10}
\end{equation}
质心速度的 $x$ 分量为
\begin{equation}
v _ {c x} = \frac {\mathrm{d} x _ {c}}{\mathrm{d} t} = \frac {1}{2} \omega l \cos \theta 
\label{eq:8.s10}
\end{equation}
质心加速度的 $x$ 分量为
\begin{equation}
a _ {c x} = \frac {\mathrm{d} v _ {c x}}{\mathrm{d} t} = - \frac {1}{2} \omega^ {2} l \sin \theta + \frac {1}{2} l \cos \theta \frac {\mathrm{d} \omega}{\mathrm{d} t}
\label{eq:9.s10}
\end{equation}
由\eqref{eq:6.s10}式得
\begin{equation}
\frac {\mathrm{d} \omega}{\mathrm{d} t} = \frac {1}{2} \omega \sin \theta \sqrt {\frac {2 g}{l (1 - \cos \theta)}} = \frac {g}{l} \sin \theta 
\label{eq:10.s10}
\end{equation}
在得到上述结果时又利用了\eqref{eq:6.s10}式。把\eqref{eq:6.s10}、\eqref{eq:10.s10}式代入\eqref{eq:9.s10}式，得
\begin{equation}
a _ {c x} = \frac {\mathrm{d} v _ {c x}}{\mathrm{d} t} = - g \sin \theta (1 - \cos \theta) + \frac {1}{2} g \sin \theta \cos \theta = g \sin \theta \left(\frac {3}{2} \cos \theta - 1\right)
\label{eq:11.s10}
\end{equation}
设竖直墙面对小球 $1$ 的正压力为 $T$，质心 $c$ 在 $x$ 方向的运动满足
\begin{equation}
T = 2 m a _ {c x}
\label{eq:12.s10}
\end{equation}
由\eqref{eq:12.s10}式可知，当 $a_{cx}=0$ 时，
\begin{equation}
T = 0,
\label{eq:13.s10}
\end{equation}
小球 $1$ 开始离开竖直墙面。由\eqref{eq:11.s10}式得，小球 $1$ 开始离开竖直墙面时，夹角 $\theta$ 满足方程
\begin{equation}
\sin \theta \left(\frac {3}{2} \cos \theta - 1\right) = 0
\label{eq:14.s10}
\end{equation}
方程\eqref{eq:13.s10}的一个解满足 $\sin\theta=0$ ，故 $\theta=0$ ，此角度对应初始位置；方程\eqref{eq:13.s10}的另一个解满足 $\frac{3}{2}\cos\theta-1=0$ ，故
\begin{equation}
\theta = \arccos \frac {2}{3}
\label{eq:15.s10}
\end{equation}
此即小球 $1$ 开始离开竖直墙面时杆与竖直墙面之间的夹角。
\end{solution}